\documentclass{article}
\usepackage[utf8]{inputenc}
\usepackage[T1]{fontenc}
\usepackage{hyperref}
\usepackage{booktabs}
\usepackage{amsmath}
\usepackage{graphicx}
\usepackage{microtype}
\usepackage{array}
\usepackage[margin=1in]{geometry}

\title{\textbf{GreenPyBench: A Benchmark for Evaluating Automated Python Code Efficiency Optimization}}

\author{
  Miguel Angel Huamani,\quad Aanya Singh Dhankhar,\quad Haoming Qin,\quad Santiago Martinez \\
  University of Illinois Urbana-Champaign \\
  CS 498 AI Agents in the Wild --- Group S11
}

\date{May 5, 2026}

\begin{document}
\maketitle

\begin{abstract}
We introduce \textbf{GreenPyBench}, a reproducible benchmark for evaluating AI agents that optimize Python programs for computational efficiency. Existing code benchmarks emphasize functional correctness but neglect runtime and memory efficiency. GreenPyBench comprises 10 procedural Python tasks spanning common inefficiency patterns: quadratic duplicate detection, linear search, loop aggregation, top-$k$ selection, nested-loop join, string concatenation, generator pipelines, memoization, data-structure choice, and multi-pass fusion. Each task includes a deterministic input generator, a slow reference implementation, an optimized reference implementation, and a PyTest suite. Agents are scored on runtime speedup, peak RAM reduction, and a composite Valid Improvement Score (VIS) that penalizes correctness regressions and code-quality degradation. Any system producing Python source code can be evaluated without modification.
\end{abstract}

%% ============================================================
\section{Introduction}
%% ============================================================

Software energy consumption at the application layer is a largely overlooked source of computational waste~\cite{strubell2019energy,schwartz2020green}. Inefficient Python scripts---featuring $O(n^2)$ loops, non-vectorized operations, redundant copies, and missed memoization---inflate runtime and peak memory in production pipelines, directly increasing energy use. Manual optimization requires profiling expertise and algorithmic reasoning that most developers lack, and existing linters rarely propose complexity-reducing changes.

Existing benchmarks (HumanEval~\cite{chen2021evaluating}, MBPP~\cite{austin2021program}, SWE-Bench~\cite{jimenez2024swebench}) measure correctness, not efficiency. PIE~\cite{shypula2023pie} targets contest-style problems and omits memory metrics and quality penalties. No benchmark provides (a) real-world procedural inefficiency patterns, (b) deterministic profiling harnesses, and (c) joint correctness-and-efficiency scoring with explicit regression penalties.

\textbf{GreenPyBench} fills this gap: it is (1) \emph{reproducible} --- deterministic seeds, fixed package versions, open harness; (2) \emph{correctness-gated} --- efficiency credit withheld for failing rewrites; (3) \emph{agent-agnostic} --- any Python-producing system can be evaluated; and (4) \emph{interpretable} --- per-task and per-family breakdowns reveal which patterns are hardest for agents.

%% ============================================================
\section{Related Work}
%% ============================================================

HumanEval~\cite{chen2021evaluating} and MBPP~\cite{austin2021program} measure function synthesis from docstrings with no efficiency metrics. SWE-Bench~\cite{jimenez2024swebench} targets repository-level bug fixing, also without profiling. PIE~\cite{shypula2023pie} pioneers efficiency as a primary metric but uses competitive-programming problems, excludes memory, and provides no quality penalties. InterCode~\cite{yang2023intercode} evaluates agents in interactive environments but focuses on correctness. GreenPyBench is the first benchmark to target the procedural inefficiency patterns common in scientific Python with a closed-loop profiling harness designed for iterative agents. It operationalizes the Green AI agenda~\cite{schwartz2020green,strubell2019energy} at the application-code layer.

%% ============================================================
\section{Benchmark Design}
%% ============================================================

\subsection{Task Domain and Scope}
GreenPyBench targets 10--50 line procedural Python programs with tractable single bottlenecks. In scope: loop-level optimizations, data-structure substitutions, $O(n^2)\!\to\!O(n)$ reductions, vectorization, and memory streaming. Out of scope: multi-file refactoring, I/O-bound tasks, GPU acceleration, and compiler-level transforms. The harness is purpose-built and does not follow an existing framework.

\subsection{Tasks}
Table~\ref{tab:tasks} summarizes all 10 tasks. Three representative tasks are described below.

\begin{table}[h]
\centering
\caption{GreenPyBench task summary. Difficulty: Easy (E), Medium (M), Hard (H).}
\label{tab:tasks}
\small
\begin{tabular}{clllc}
\toprule
\textbf{\#} & \textbf{Task} & \textbf{Inefficiency} & \textbf{Optimization} & \textbf{Diff.} \\
\midrule
1  & Duplicate Detection   & $O(n^2)$ nested loop       & Hash-based lookup      & M \\
2  & Linear Search         & Full-list scan             & Set membership/bisect  & E \\
3  & Loop Aggregation      & Python loop over array     & NumPy vectorization    & E \\
4  & Top-K Selection       & Full sort for partial res. & Heap-based selection   & M \\
5  & Nested-Loop Join      & $O(n \cdot m)$ cross-join  & Hash join              & H \\
6  & String Concatenation  & Repeated \texttt{+} in loop & \texttt{str.join}     & E \\
7  & Generator Pipeline    & Intermediate list alloc.   & Generator chaining     & M \\
8  & Memoization           & Redundant recursive calls  & LRU cache / DP         & M \\
9  & Data Structure Choice & List where set fits        & Set/dict substitution  & M \\
10 & Multi-Pass Fusion     & Three separate traversals  & Single-pass fusion     & H \\
\bottomrule
\end{tabular}
\end{table}

\paragraph{Duplicate Detection (Medium).} Input: $n \in \{10^4,10^5,10^6\}$ integers. Slow: $O(n^2)$ nested comparison. Goal: single-pass frequency table yielding $O(n)$. Output normalized by sorting to ensure test independence.

\paragraph{Nested-Loop Join (Hard).} Input: two tables of $n\!=\!m\!=\!2{,}000$ dicts with a shared key. Slow: $O(n \cdot m)$ double loop. Goal: hash-index one table for $O(n+m)$ join. Output sorted by key before comparison.

\paragraph{Multi-Pass Fusion (Hard).} Input: $n \in \{10^5,10^6\}$ floats. Slow: three separate passes computing sum, mean, std with intermediate allocations. Goal: single-pass running-sum-of-squares fusion. Output matched to six decimal places.

\subsection{Design Principles and Difficulty}
Tasks span five pattern families: algorithmic reduction (1,4,5), data-structure substitution (2,9), vectorization (3), memory streaming (7,10), and caching (8). All slow implementations are idiomatic Python a competent developer might write, not adversarially obfuscated code. Correctness tests normalize outputs before comparison. Table~\ref{tab:difficulty} shows the difficulty distribution; harder tasks offer larger optimization headroom, validating the labels.

\begin{table}[h]
\centering
\caption{Difficulty distribution. Avg.\ speedup is reference optimized vs.\ slow.}
\label{tab:difficulty}
\small
\begin{tabular}{lccl}
\toprule
\textbf{Level} & \textbf{Count} & \textbf{Avg.\ Speedup} & \textbf{Tasks} \\
\midrule
Easy   & 3 & 4.1$\times$  & 2, 3, 6 \\
Medium & 5 & 11.3$\times$ & 1, 4, 7, 8, 9 \\
Hard   & 2 & 38.7$\times$ & 5, 10 \\
\bottomrule
\end{tabular}
\end{table}

%% ============================================================
\section{Evaluation Protocol}
%% ============================================================

\paragraph{Metrics.} Per task: (1) \emph{Correctness pass rate} over PyTest suite; (2) \emph{Runtime speedup} (\%) --- median of 7 measured runs (after 2 warm-ups) in an isolated subprocess; (3) \emph{Peak RAM reduction} (\%) --- median \texttt{tracemalloc} snapshot.

\paragraph{Scoring.} Correctness is a hard gate. The per-task Valid Improvement Score (VIS) is:
\begin{equation}
\text{VIS} = \max\!\left(0,\; 0.7\times\text{runtime\_gain} + 0.3\times\text{memory\_gain} - \text{quality\_penalty}\right)
\end{equation}
Gains are clipped at zero. $\text{quality\_penalty}=0.25$ if cyclomatic complexity increases $>$20\% or lines of code increase $>$30\% without achieving $\geq$10\% speedup or $\geq$5\% RAM reduction (measured via \texttt{radon}).

\paragraph{Aggregation and Reporting.} Report mean correctness, mean runtime speedup, mean RAM reduction, mean VIS, and regression rate---broken down by difficulty tier and pattern family. LLM-based methods should be run at least three times to estimate variance. Submissions must specify model name, temperature, iteration limit, and Python environment.

\paragraph{Model declaration for this submission.} The reported one-trial LLM numbers were generated with Tinker using \texttt{Qwen/Qwen3-30B-A3B-Instruct-2507} as the shared model across Single-Shot, One-Pass, and Agent settings.

%% ============================================================
\section{Benchmark Validation}
%% ============================================================

\paragraph{Harness reliability.} Running each slow implementation 20 times yielded CV $<5\%$ for runtime and CV $<3\%$ for RAM across all tasks, confirming that 7-run median aggregation produces stable estimates.

\paragraph{Baseline results.} Table~\ref{tab:baseline} reports four conditions using one-trial evaluation. The static rule baseline applies deterministic rewrites (\texttt{join}, \texttt{@lru\_cache}, NumPy, \texttt{set()}, \texttt{heapq}) and returns the original when no rule fires; it remains high-correctness but can regress on specific tasks. Compared with single-shot prompting, one-pass profile guidance improves robustness but is sensitive to outliers (notably Multi-Pass Reduction). These outcomes support the benchmark's goal of evaluating both optimization quality and stability.

\begin{table}[h]
\centering
\caption{Baseline results on GreenPyBench (10 tasks, one-trial). RT = runtime speedup, RAM = peak RAM reduction, both means over passing tasks.}
\label{tab:baseline}
\small
\begin{tabular}{lcccc}
\toprule
\textbf{System} & \textbf{Correct (\%)} & \textbf{RT Speedup (\%)} & \textbf{RAM Red.\ (\%)} & \textbf{VIS} \\
\midrule
Original slow code        & 100 & 0.0  & 0.0  & 0.00 \\
Static rule baseline      & 100 & 25.6 & 0.1  & 18.32 \\
Single-shot LLM           & 10 & 19.9  & 0.0  & 13.92 \\
One-pass profile-guided   & 90 & -45.8 & 12.5 & 45.98 \\
Reference optimized impl. & 100 & 82.1 & 41.6 & --- \\
\bottomrule
\end{tabular}
\end{table}

\paragraph{Human baseline.}
\begin{table}[h]
\centering
\caption{Manual baseline on three representative tasks (7-run median harness evaluation).}
\small
\begin{tabular}{lccc}
\toprule
Task & Speedup (\%) & Correct (\%) & Time (min) \\
\midrule
Duplicate Detection & 99.9 & 100 & 6 \\
Nested-Loop Join    & 94.9 & 100 & 8 \\
Memoization         & 100.0 & 100 & 6 \\
\bottomrule
\end{tabular}
\end{table}
The manual baseline confirms that straightforward algorithmic rewrites already deliver substantial gains on representative tasks, while still requiring explicit implementation effort and validation via the same harness protocol.

%% ============================================================
\section{Conclusion}
%% ============================================================

GreenPyBench is a reproducible benchmark for AI agents on Python code efficiency, with joint correctness-and-efficiency metrics, regression penalties, and a profiling harness for iterative systems. The one-trial baselines already show that static rules and single-shot prompting are insufficiently robust across tasks, while profiler-guided and iterative methods achieve stronger average VIS. The implementation supports Tinker-compatible API routing for budget-aware reproducibility. The harness, tasks, seeds, and baselines will be released to support reproducible Green AI research. Future extensions include multi-trial confidence intervals, energy measurement via RAPL counters, and evaluation on real-world scientific codebases.

\paragraph{Changelog.} Final version adds confirmed per-task static rule VIS results, refines \texttt{tracemalloc} harness details, expands task design principles, includes a manual baseline on Tasks 1, 5, and 8, and updates evaluation infrastructure notes for Anthropic/Tinker endpoint compatibility.

%% ============================================================
\begin{thebibliography}{10}
\bibitem{austin2021program} J. Austin et al. Program synthesis with large language models. \emph{arXiv:2108.07732}, 2021.
\bibitem{chen2021evaluating} M. Chen et al. Evaluating large language models trained on code. \emph{arXiv:2107.03374}, 2021.
\bibitem{jimenez2024swebench} C.E. Jimenez et al. SWE-bench: Can language models resolve real-world GitHub issues? \emph{arXiv:2310.06770}, 2024.
\bibitem{madaan2023chatgpt} A. Madaan et al. Is your code generated by ChatGPT really correct? \emph{arXiv:2305.01210}, 2023.
\bibitem{schwartz2020green} R. Schwartz et al. Green AI. \emph{Commun.\ ACM}, 63(12):54--63, 2020.
\bibitem{shypula2023pie} A. Shypula et al. Learning performance-improving code edits. \emph{arXiv:2302.07867}, 2023.
\bibitem{strubell2019energy} E. Strubell et al. Energy and policy considerations for deep learning in NLP. \emph{arXiv:1906.02629}, 2019.
\bibitem{xia2023evoeval} C.S. Xia et al. EvoEval: Evolving coding benchmarks via LLMs. \emph{arXiv:2403.19114}, 2023.
\bibitem{yang2023intercode} J. Yang et al. InterCode: Standardizing and benchmarking interactive coding. \emph{arXiv:2306.14898}, 2023.
\end{thebibliography}

\end{document}
